\documentclass[a4paper, 12pt]{article}

\usepackage[utf8]{inputenc}
\usepackage[T2A]{fontenc}
\usepackage[english,russian]{babel}
\usepackage{amsmath, amsfonts, amsthm, amssymb, mathtools}
\usepackage{multicol}
\usepackage{enumitem}
\usepackage{cancel}
\usepackage{hyperref}
\hypersetup{
    colorlinks,
    linkcolor={black},
    citecolor={black},
    urlcolor={blue},
}
\usepackage{indentfirst}

\newcommand{\T}{\mathcal{T}}
\newcommand{\w}{\omega}
\newcommand{\rev}{\operatorname{rev}}

\title{Лабораторная работа №1}
\author{Чайкин Семён\\Вариант 26}

\begin{document}
    \maketitle
    \tableofcontents
    \newpage

    \section{Задание}    

    Дана SRS $\T$:

    \begin{gather*}
        ab  \to br\\
        raa \to dr\\
        ara \to ad\\
        da  \to aad\\
        dd  \to abdar 
    \end{gather*}

    По имеющейся SRS определить:

    \begin{itemize}
        \itemsep -0.8pt
        \item[---] завершимость;
        \item[---] конечность классов эквивалентности по НФ;
        \item[---] локальную конфлюэнтность и пополняемость по Кнуту-Бендиксу.
    \end{itemize}

    \section{Исследование SRS}

    \subsection{Завершимость}

    \textbf{Без понятия}.
    
    \subsection{Классы эквивалентности}

    Рассмотрим слова $b^n$ ($n \geq 1$) --- каждое из них является НФ. Очевидно, что каждое подобное слово лежит в собственном классе эквивалентности, так как при любой направлении стрелок в SRS, не существует правил для переписывания строк состоящих только из символа $b$. Таким образом, в $\T$ \textbf{бесконечное} количество классов эквивалентности.

    \subsection{Локальная конфлюэнтность и пополнимость}

    \subsubsection{Переупорядочивание правил}

    Переориентируем правила, чтобы задать фундированный порядок. Рассмотрим, например, следующую SRS:

    \begin{gather*}
    br    \to ab\\
    raa   \to dr\\
    ara   \to ad\\
    aad   \to da\\
    abdar \to dd
    \end{gather*}

    Эта SRS завершима, так как каждое правило уменьшает длину слова.
    
    \textbf{Фундированный порядок}: length-lexicographic (\textit{армейский}), где $a < b < d < r$.

    \subsubsection{Проверка локальной конфлюэнтности}

    Рассмотрим слово $braa$:

    \begin{gather*}
        \underset{ab}{\underline{br}} aa \to abaa\\
        b\underset{dr}{\underline{raa}} \to bdr
    \end{gather*}

    Получившиеся справа слова являются НФ. Как мы видим, SRS \textbf{не является} локально конфлюэнтной.

    \subsubsection{Пополняемость по Кнуту-Бендиксу}

    Рассмотрим два правила из $\T$: $raa \to dr$ и $aad \to da$ --- они формируют критическую пару $\langle \mathbf{dr}ad, ra\mathbf{da} \rangle$, получаемую переписыванием строки $raaad$. Обе части пары нельзя редуцировать дальше, поэтому добавляем в систему новое правило $rada \to drad$.

    Рассмотрим теперь правила $ara \to ad$ и $rada \to drad$. Они тоже формируют критическую пару $\langle \mathbf{ad}da, a\mathbf{drad} \rangle$, получаемую переписыванием строки $arada$. Обе части пары нельзя редуцировать дальше, поэтому добавляем в систему новое правило $adrad \to adda$.

    Рассмотрим правила $rada \to drad$ и $adrad \to adda$. Они формируют критическую пару $\langle ad\mathbf{drad}, \mathbf{addaa} \rangle$, получаемую переписыванием строки $adrada$. Обе части пары нельзя редуцировать дальше, поэтому добавляем новое правило $addrad \to addaa$.

    Заметим, что если повторно рассмотреть только что получившееся правило $ad^2rad \to adda^2$ с правилом $rada \to drad$, мы получим новое правило $ad^3 rad \to adda^3$. Продолжая эту процедуру для правил $rada \to drad$ и $ad^nrad \to adda^n$ ($n \geq 1$), получаем новое правило $ad^{n+1}rad \to adda^{n+1}$, которое также нельзя редуцировать, поэтому систему $\tilde{\T}$ можно пополнять до бесконечности. Таким образом, система является непополнимой по Кнуту-Бендиксу при указанной выше стратегии выбора критических пар. 

    При этом из работы \href{https://github.com/ElizaaaSh/TFL}{@ElizaaaSh} видно, что система является пополнимой при выборе другой стратегии.

    \subsection{Построение $\T'$}

    На основе исследования SRS $\tilde{\T}$ была построена эквивалентная ей SRS $\T'$:

    \begin{flushleft}
    \begin{multicols}{4}
    $br      \to ab$\\
    $raa     \to dr$\\
    $ara     \to ad$\\
    $aad     \to da$\\
    $abdar   \to dd$\\
    $abdad   \to dda$\\
    $ddar    \to ddaa$\\
    $adbdar  \to ardd$\\
    $arda    \to adad$\\
    $adra    \to add$\\
    $adr     \to ada$\\
    $drbdar  \to radd$\\
    $rda     \to drd$\\
    $rada    \to drad$\\
    $drra    \to drd$\\
    $abaa    \to bdr$\\
    $adrad   \to adda$\\
    $addrad  \to addaa$\\
    $adddrad \to addaaa$\\
    \end{multicols}
    \end{flushleft}

    \subsection{Инварианты}

    Исходная SRS $\T$ (после изменения порядка):

    \begin{gather*}
    br    \to ab\\
    raa   \to dr\\
    ara   \to ad\\
    aad   \to da\\
    abdar \to dd
    \end{gather*}

    Итоговая SRS $\T'$:

    \begin{flushleft}
    \begin{multicols}{4}
    $br      \to ab$\\
    $raa     \to dr$\\
    $ara     \to ad$\\
    $aad     \to da$\\
    $abdar   \to dd$\\
    $abdad   \to dda$\\
    $ddar    \to ddaa$\\
    $adbdar  \to ardd$\\
    $arda    \to adad$\\
    $adra    \to add$\\
    $adr     \to ada$\\
    $drbdar  \to radd$\\
    $rda     \to drd$\\
    $rada    \to drad$\\
    $drra    \to drd$\\
    $abaa    \to bdr$\\
    $adrad   \to adda$\\
    $addrad  \to addaa$\\
    $adddrad \to addaaa$\\
    \end{multicols}
    \end{flushleft}


    \textbf{Инварианты}:
    \begin{enumerate}
        \itemsep -0.8pt 
        \item Количество букв $b$ не возрастает.
        \item Количество букв $r$ не возрастает.
        \item Количество букв $d$ не убывает.
        \item Если слово заканчивается на букву $b$, то после любой последовательности преобразований, оно будет заканчиваться на $b$.
        \item Если слово начинается на букву $d$, то после любой последовательности преобразований, оно будет начинаться на $d$.
    \end{enumerate}
    
\end{document}

